\subsection{Microservicio de Laboratorios y Recursos}
\label{subsec:laboratorios-recursos}

\subsubsection{Alcance y responsabilidad}

El microservicio mantiene el catálogo de laboratorios y equipos de CLAB,
incluidas las condiciones operativas que restringen su uso. Su motivo de
cambio son las reglas de administración de los recursos físicos. El alcance
se definió a partir del código del sistema y de las historias de usuario de
la Entrega~1, siguiendo el criterio de analizar el dominio y sus dependencias
para identificar los límites del servicio
\cite{abgaz2023decomposition,oumoussa2024evolution}.

El servicio gestiona laboratorios y su equipamiento (RF-03), el inventario
por número de serie (RF-07), el historial de estado de equipos (RF-08) y
el registro de bloqueos institucionales (RF-11). También administra las
asignaciones de encargados y las referencias a fotografías. Para el RF-04,
aporta las condiciones operativas y los bloqueos vigentes. Reservas y Agenda
combina esa información con su calendario y decide qué franjas puede
ofrecer o confirmar. El módulo de reportes registra y sigue las fallas
(RF-09 y RF-10) usando los identificadores de laboratorio y equipo.

\subsubsection{Propiedad de datos}

El servicio usa una base PostgreSQL propia, con las siguientes entidades:

\begin{itemize}
    \item \textbf{Laboratorio.} Conserva ubicación, capacidad, descripción,
    estado operativo e identificador de sede. Su equipamiento se consulta
    a partir de los registros del inventario.
    \item \textbf{Equipo y tipo de equipo.} Cada equipo pertenece a un
    laboratorio y a un tipo. Registra número de serie, marca, modelo y estado.
    El número de serie es obligatorio y único en la base.
    \item \textbf{Historial de estado.} Cada cambio a operativo, mantenimiento
    o fuera de servicio conserva el equipo, el estado anterior y el nuevo,
    la fecha y el identificador del usuario responsable. El cambio de estado
    y su historial se guardan en una sola transacción.
    \item \textbf{Bloqueo y tipo de bloqueo.} Cada bloqueo identifica el
    laboratorio, el tipo, el motivo, el intervalo y el estado; también
    conserva el identificador del usuario que lo registró.
    \item \textbf{Asignación de encargado y referencia de fotografía.}
    La asignación vincula al usuario con el laboratorio y conserva su
    vigencia. Las referencias de fotografía identifican la imagen asociada
    a un laboratorio o equipo.
    \item \textbf{Evento pendiente.} Conserva los cambios de bloqueos que
    aún deben publicarse en RabbitMQ.
\end{itemize}

PostgreSQL conserva el motor utilizado por CLAB y permite aplicar
unicidad, claves foráneas internas y transacciones locales. Los usuarios
y las sedes se referencian mediante identificadores externos. Las reservas,
su ocupación y las acciones sobre reservas afectadas pertenecen a
Reservas y Agenda.

\subsubsection{Componentes internos}

\begin{description}
    \item[API REST.] Recibe las solicitudes del Gateway y la consulta
    interna de Reservas y Agenda. Valida la entrada y la dirige al componente
    correspondiente.
    \item[Catálogo y encargados.] Gestiona laboratorios y asignaciones.
    Consulta al servicio de identidad para validar al encargado.
    \item[Inventario e historial.] Gestiona equipos y tipos de equipo,
    y registra los cambios de estado con fecha y usuario.
    \item[Gestión de bloqueos.] Registra, modifica y levanta bloqueos.
    Persiste el cambio y el evento pendiente en la misma transacción.
    \item[Consulta de condiciones operativas.] Consulta el estado del
    laboratorio y los bloqueos vigentes que se solapan con el intervalo.
    \item[Referencias de fotografías.] Mantiene los vínculos entre las
    imágenes y sus laboratorios o equipos.
    \item[Repositorios de persistencia.] Leen y escriben los datos de la
    base PostgreSQL del servicio.
    \item[Publicador de eventos.] Recupera eventos pendientes, los envía
    a RabbitMQ y registra su publicación al recibir la confirmación del
    intermediario. Reintenta los envíos sin confirmar.
\end{description}

\subsubsection{Mecanismos de comunicación}

\paragraph{REST sobre HTTPS, sincrónico.}
Los clientes web y móvil ingresan por el API Gateway, que autentica,
enruta y limita las solicitudes. Reservas y Agenda consulta el estado y
los bloqueos vigentes del laboratorio antes de ofrecer o confirmar una
franja. Si falla esa consulta, pospone la confirmación. La validación de
un encargado también usa REST porque la asignación necesita el resultado
de identidad antes de guardarse.

\begin{description}
    \item[Catálogo.] \texttt{GET /laboratorios},
    \texttt{POST /laboratorios} y \texttt{PATCH /laboratorios/\{id\}}.
    \item[Inventario.] \texttt{GET /equipos}, \texttt{POST /equipos} y
    \texttt{PATCH /equipos/\{id\}/estado}. Un número de serie duplicado
    provoca el rechazo del registro.
    \item[Bloqueos.] \texttt{GET /bloqueos}, \texttt{POST /bloqueos} y
    \texttt{PATCH /bloqueos/\{id\}}. La última operación permite
    modificar o levantar un bloqueo.
    \item[Consulta interna.] \texttt{GET /laboratorios/\{id\}/}\allowbreak
    \texttt{condiciones-operativas}. Recibe el intervalo solicitado y
    devuelve el estado operativo y los bloqueos vigentes aplicables.
\end{description}

\paragraph{RabbitMQ sobre AMQP, asincrónico.}
El servicio publica los eventos \texttt{BloqueoLaboratorioRegistrado},
\texttt{BloqueoLaboratorioModificado} y
\texttt{BloqueoLaboratorioLevantado}. Cada evento incluye su identificador,
los identificadores del laboratorio y del bloqueo, el intervalo, el estado
y la versión del bloqueo. Reservas y Agenda usa estos avisos para reevaluar
las reservas afectadas. Registra los identificadores ya procesados y
descarta versiones anteriores del mismo bloqueo.

REST permite consultar las restricciones vigentes en el momento de decidir.
RabbitMQ permite comunicar los cambios sin esperar a que Agenda los procese
y el registro local permite reintentar si el intermediario se interrumpe.
La consulta REST y la escritura de la reserva ocurren en servicios
distintos. Si se registra un bloqueo entre ambas operaciones, el evento
posterior activa la reevaluación de la reserva. Estos mecanismos se
justifican como criterios de diseño a partir de los flujos de la Entrega~1.
